% Copyright 2026 The terCAD team. All rights reserved.
% Use of this content is governed by a CC BY-NC-ND 4.0 license that can be found in the LICENSE file.

\subsection{Why Set Goals}
\markboth{Planning}{Why Set Goals}

A budget answers the question 'Where does the money go?', while a goal answers 'What is it
all for?'. Without a clear destination, saving can feel like a sacrifice with no reward, and
any surplus is quietly absorbed by everyday spending. However, with a goal, every unspent
amount becomes a visible step towards something meaningful: a holiday, a new laptop, a
deposit on a property, or simply peace of mind.

A widely used checklist for setting well-formed goals is known by the acronym \emph{SMART}.
A goal should be:

\begin{itemize}
  \item \textbf{Specific} -- \q{a vacation in Portugal}, not \q{to travel more};
  \item \textbf{Measurable} -- with a target amount, so that progress can be tracked;
  \item \textbf{Achievable} -- realistic with respect to the current income and expenses;
  \item \textbf{Relevant} -- important enough to justify the effort;
  \item \textbf{Time-bound} -- with a deadline that creates a pace.
\end{itemize}

\noindent Fingrom addresses the measurable and time-bound aspects directly: each goal has a
target amount and a deadline, and the app calculates how much has been saved towards the goal
and how much needs to be saved each month to reach it on time (\cref{img:goal-view}).
However, deciding what is specific, achievable and relevant remains a personal decision.

\img{img/ch04-s01-goal__view}{Goal View}{img:goal-view}


\subsubsection{Goals and Budgets}

Goals and budgets are two sides of the same plan. Budget categories limit expenses, while
goals determine how the difference between income and expenditure is allocated. The
well-known principle 'pay yourself first' prioritises goals: first decide how much to set
aside, then distribute the rest among the budget categories. For example, the 50/30/20 rule
(see \cref{budget:limit}) allocates 20\% to savings, which is exactly what is needed to
achieve your goals.

\subsubsection{Types of Goals}

Goals differ mainly by their horizon:

\begin{center}
\small
\begin{tabularx}{\textwidth}{|l|l|X|}
  \hline
  \textbf{Horizon} & \textbf{Period} & \textbf{Examples} \\
  \hline
  Short-term & up to 1 year & A new phone, a vacation, holiday gifts, an annual insurance
  payment. \\
  \hline
  Mid-term & 1 -- 5 years & A car, a wedding, renovation, an emergency fund of 3 -- 6 months
  of expenses. \\
  \hline
  Long-term & over 5 years & A down payment for a home, education of children, early
  retirement. \\
  \hline
\end{tabularx}
\end{center}

\noindent An \emph{emergency fund} deserves a special mention. It is usually recommended as
the first goal, as having a reserve that covers several months of essential expenses protects
all your other goals (and your budget) from being derailed by unexpected events such as job
loss, illness or urgent repairs.


\subsection{Creating a Goal}
\markboth{Planning}{Creating a Goal}

Goals are managed on the \textbf{Goals} page, where a new one is added via the
\textbf{Add new Goal} button. The form contains the following fields:

\begin{description}
  \item[Goal] (required) -- a short description, e.g. \q{Vacation in Portugal}.
  \item[Target Amount] -- the amount to be collected.
  \item[Currency] -- the currency of the target amount. It may differ from the default
  currency (e.g. a vacation planned in euros while the income is in dollars); the progress is
  calculated using the exchange rates.
  \item[Finished till Date] -- the deadline.
  \item[Icon] and \textbf{Color} -- visual identification of the goal in lists.
\end{description}

\noindent When the goal is created, Fingrom memorizes the current total of all accounts as
the starting point of the goal. Money that already exists is not counted towards a new goal;
only money saved from that point onwards moves the progress bar. This ensures that goals are
realistic -- setting a goal of £1,000 while having £5,000 in your accounts does not instantly
complete it.


\subsection{How Progress Is Calculated}
\label{goal:progress}
\markboth{Planning}{How Progress Is Calculated}

Fingrom does not require money to be transferred to a separate goal account. Instead, goal
progress follows the total net worth: whenever it grows, the goals move forward; whenever it
shrinks, they move back.

\subsubsection{Distribution of Changes}

Let $A$ be the target amount of a goal and $p$ its current progress (from $0$ up to $1$).
When the total of accounts changes by $\Delta$ -- after an income, a bill, or an edit of a
balance -- the change is distributed equally among all $N$ active goals, and the progress of
each of them is updated as:
\[
  p' = \frac{A \cdot p + \Delta / N}{A}
\]

\noindent with the following rules:

\begin{itemize}
  \item the progress never falls below $0$ and never exceeds $1$ (100\%);
  \item when a goal reaches 100\%, the excess of its share is passed on to the remaining
  goals, so no saved money is lost;
  \item while the total stays below the starting point of a goal, growth only restores the
  losses and does not count as progress;
  \item a completed goal remains completed as long as the total of accounts still covers it,
  even if a part of the money is spent afterward;
  \item amounts are converted between the default currency and the currency of each goal.
\end{itemize}

\noindent Transfers between own accounts do not change the total, and therefore do not
influence goals at all.

\subsubsection{Example}

Suppose there are £5,000 in the accounts and two goals are created: 'Vacation' with a target
of £1,200 and 'Laptop' with a target of £1,500.

\begin{center}
\footnotesize
\begin{tabular}{|l|r|r|r|}
  \hline
  \textbf{Event} & \textbf{Total} & \textbf{Vacation} & \textbf{Laptop} \\
  \hline
  Goals created & 5000 & 0\% (0) & 0\% (0) \\
  Salary \q{+2000}, bills \q{-1400} & 5600 & 25\% (300) & 20\% (300) \\
  A gift \q{+200} & 5800 & 33\% (400) & 27\% (400) \\
  Unexpected repair \q{-300} & 5500 & 21\% (250) & 17\% (250) \\
  \hline
\end{tabular}
\end{center}

\noindent Each change is split between the two goals. As their targets differ, the same
amount represents a different percentage. For example, 300 is a quarter of the holiday
budget, but only a fifth of the budget for the laptop.

\subsubsection{Time Progress}

Besides the collected amount, each goal has a \emph{time progress} -- the share of time
already passed between its creation and its deadline:
\[
  t = \frac{\textit{days since creation}}{\textit{days from creation till deadline}}
\]

\noindent Comparing both values tells whether the goal is on schedule. If $p \geq t$ (e.g.
40\% collected after 30\% of time), the goal is ahead; if $p < t$, saving has to speed up --
or the deadline has to be moved.

\subsubsection{Editing a Goal}

The target amount, the deadline, and other fields can be changed at any time. The already
collected amount is preserved, and the progress is recalculated against the new target:
\[
  p' = \frac{A_{old} \cdot p}{A_{new}}
\]

\noindent For example, if the target is set at 1,200 and 50\% of this is collected, raising
the target to 1,500 would mean that the same 600 would represent 40\%.


\subsection{Monthly Target}
\markboth{Planning}{Monthly Target}

While knowing the progress is useful, the practical question is 'How much should I save this
month?'. The \textbf{Profit vs Goals Target} block at the top of the \textbf{Goals} page
(\cref{img:goal-view}) provides the answer. For each active goal, the remaining amount is
divided by the number of months until the deadline (at least one) and the results are added
together:
\[
  \textit{Goals Target} = \sum_{goals} \frac{A \cdot (1 - p)}{\max(1, \textit{months left})}
\]

\noindent This target is compared with the current month's actual result, both expressed in
the default currency.

\begin{itemize}
  \item \textbf{Invoices} -- the incomes of the current month (transfers are not counted);
  \item \textbf{Bills} -- the expenses of the current month;
  \item \textbf{Profit} -- incomes minus bills, i.e. what has actually been saved.
\end{itemize}

\noindent A gauge illustrates the comparison: the monthly target is positioned slightly
before the middle of the scale. Therefore, a needle beyond the middle indicates that the
month will exceed the required goals, while a needle on the left signals a shortfall.

For example, 'Vacation' needs £900 more within six months (£150 per month), and 'Laptop'
needs £1,250 more within five months (£250 per month). The target goal is 400. With an income
of £2,000 and bills of £1,500 in the current month, the profit of £500 exceeds the target --
the plan is on track.

Please note that if a deadline has already passed, one month is added to the target for the
current month. This is a signal to either complete the goal or move its deadline.


\subsection{Completing a Goal}
\markboth{Planning}{Completing a Goal}

Once a goal has reached 100\%, it is marked as achieved on both the main page and the
\textbf{Goals} page. The \textbf{Complete Goal} action then becomes available in its details.
Completing a goal means spending the collected money on its intended purpose:

\begin{enumerate}
  \item a new bill for the target amount is created, titled \q{Complete Goal: \dots} with the
  name of the goal;
  \item the bill opens for editing -- select the account the money is paid from and the
  budget category, and adjust the amount and date to the actual purchase;
  \item the goal itself is deactivated and no longer takes part in the distribution of
  savings.
\end{enumerate}

\noindent  Any goal that is no longer relevant can be deleted via the \textbf{Delete Goal}
action. The progress of the remaining goals is unaffected in both cases, and each of them
receives a larger share of future savings from that moment onwards.


\subsection{Goals on the Main Page}
\markboth{Planning}{Goals on the Main Page}

Active goals can be displayed on the main page with a progress bar, showing the collected and
target amounts and the deadline. The header of the block shows the total amount of all the
goals -- the sum of the remaining amounts to be collected, converted into the default
currency:
\[
  \textit{Goals Total} = \sum_{goals} A \cdot (1 - p)
\]

\noindent Like other components, the goals block can be added, moved, or removed via
\textbf{Customize Page}.


\subsection{Practical Tips}
\markboth{Planning}{Practical Tips}

\textbf{Start with an emergency fund.} It protects all other goals from being consumed by the
first unexpected expense.

\textbf{Keep the number of goals small.} Savings are split equally among active goals; with
ten goals at once, each of them moves slowly. Two or three goals at a time are easier to
achieve.

\textbf{Set realistic deadlines.} Check the \textbf{Goals Target} right after creating a
goal: if it exceeds the usual monthly profit, the deadline is too close or the target too
high.

\textbf{Compare money and time.} A goal that is collected slower than the time passes needs
attention now, not a week before the deadline.

\textbf{Review goals monthly.} Together with the budget review, check the progress, move
deadlines if needed, and complete or delete goals that are achieved or no longer relevant.

\textbf{Celebrate.} Completing a goal is the reward for months of discipline -- record the
purchase via \textbf{Complete Goal} and set the next one.
